无限长单原子链中除 $n=0$ 处的质量为 $M$ 外，所有格点质量均为 $m>0$；每对相邻粒子由劲度系数为 $k>0$ 的相同轻弹簧连接，平衡间距为 $a>0$，忽略耗散。取缺陷比主体质量轻，即 $0<M<m$。
(1) 求无缺陷时的行波色散关系和声学频带上边缘；(2) 设缺陷诱导的局域振动形式为 $u_n(t)=A(-1)^n e^{-\lambda|n|}e^{-i\omega t}$，其中 $\lambda>0$，由非缺陷格点方程求 $\omega(\lambda)$；(3) 使用缺陷点方程解出 $\lambda$ 和 $\omega$ 的精确表达式，证明频率位于正常频带之外；(4) 研究 $M\to m^-$ 和 $M\to0^+$ 的局域长度与频率极限，解释重质量缺陷是否能以同样形式产生高频局域模。
